\section{The Moral Universe}

\begin{quotex}

First and foremost, I will tell you who should work in this work, and when, and by what means: and what discretion you shall have in it. If you ask me who shall work thus, I answer you—all who have forsaken the world in a \textbf{true will}. \flright{\emph{The Cloud of Unknowing}}

\end{quotex}
If you start at the Greek and Roman antiquities section of the Louvre, then the transition to the Medieval art exhibit is striking. So much so, that it does not seem possible that they could have been created by the same peoples. In the same way, the transition from the Naturalistic poetry of Lucretius to the Moral Universe of Dante also cries out for an explanation. Keep in mind that Virgil, who aided Dante on his journey, had studied, like Lucretius, with the Epicureans.

\paragraph{Naturalism}
For naturalism, the purposeless movement of atoms leaves no space for moral acts. It recognizes a single force, the élan vital or vital impulse. The moral universe requires a countervailing force, the frein vital, which is the vital control that acts as a check on the impulse. That naturalistic goal is the satisfaction of the appetites, and any check on the appetite comes from an external source, not from an inner check.

\paragraph{Rationalism}
\textbf{George Santayana} in \emph{Three Philosophical Poets} asserts that Socrates provides the link between Lucretius and Dante. In the Phaedo, while sitting in prison, Socrates recounts how in his youth he had sought the ultimate cause of the Universe. The Epicureans claimed it was atoms. Other thinkers had proposed alternatives such as Water, Fire, Air. Socrates rejected them all, favoring instead Reason or Mind as the source of order. In that case, things are to be understood by their use or purpose, and not by the initial element that set them in motion.

This is the beginning of Absolute Idealism, since Being is known by Thought and not by the physics behind it. In this case, Being is the object of contemplation. This is cataphatic philosophy which seeks to make positive statements about the Absolute. The three main strands among the Greeks offer different qualities as the object of contemplation:

\begin{itemize}
\item \textbf{Aristotle}: The Unmoved Mover 
\item \textbf{Plato}: The Good 
\item \textbf{The Stoics}: The Logos 
\end{itemize}
Still, in these formulations, the Absolute is passive in relationship to Man, who is a Subject, an I. With its historical revelation, Christianity went beyond that. It was revealed that the Absolute itself was also a Subject. Obviously, the Absolute as Subject is the God, now an active force. The three qualities now take on a deeper dimension.

\begin{itemize}
\item \textbf{The Unmoved Mover}. The unmoved mover that sets the universe in motion becomes, as Subject, the \textbf{Creator} of the universe. 
\item \textbf{The Good}. The Absolute as subject seeks, or wills, the Good. That is, the Absolute is \textbf{Love}. 
\item \textbf{The Logos}. The order of the Universe is not merely the object of contemplation, but as subject it is the \textbf{Christ}. Any disorder in the world must be rectified by a Savior incarnating into the World Process. 
\end{itemize}
\paragraph{The World as Magic}
Santayana raises the question of how does Love act on the physical world. Science is insufficient to explain the World. A moral interpretation of Nature is necessary. Santayana describes the action of the moral universe on the physical as magical. He writes:

\begin{quotex}
[The Absolute] works only in its capacity of ideal; therefore, even if it exists, it works only by magic. The matter beneath feels the spell of its presence, and catches something of its image, as the waves of the sea might receive and reflect tremblingly the light shed by the moon. The world accordingly is moved and vivified in very fibre by magic, by the magic of the goal to which it aspires. 

\end{quotex}
Santayana explains that on Earth, the Magic is called Love. Actual things are symbols for what they ultimately should be. That is, existing phenomenal things are a theophany of something higher, and it requires the powers of a spiritual vision to discern that. Santayana describes it:

\begin{quotex}
We should attribute the power which things exerted over us, not to the rarer or denser substance, but to the eternal ideas that they existed by expressing, and existed to express. Things merely localized—like the saint's relics—the influences which flowed to us from above. In the world of values they were mere symbols, accidental channels for divine energy; and since divine energy, by its magic assimilation of nature, had created these things, in order to express itself, they were symbols altogether not merely in their use, but in their origin and nature. 

\end{quotex}
So for Dante, the world is love, magic, and symbolism. The pessimist Santayana has this vain hope:

\begin{quotex}
A thousand years after Dante, we may hope that his conscientious vision of the universe, where all is \emph{love, magic, and symbolism}, may charm mankind exclusively as poetry. 

\end{quotex}
Of course, that is not sufficiently pessimistic, since in our time Dante is not even appreciated as poetry. The lure of naturalism, the worldview of Satan, persists in obscuring love, magic, and symbolism from the World. For the naturalist, the solution to the mystery of life is to maintain adequate levels of serotonin, dopamine, and oxytocin, pharmacologically if necessary.

How radically different is the world filled with love, magic, and symbols of spiritual reality from the blind movement of atoms. No, Dante is not merely a poet, he is a guide showing us the way from Hell to Paradise.

\paragraph{Demanding a Choice}
Curiously, Santayana has the incisive intelligence to discern what is really at stake between the naturalistic and moral views of the world, even if he has personally chosen the path of naturalism. Paradoxically, for the naturalist, there is no choice; the mind is an illusion and everything is determined by the motion of atoms. For the moralist, on the other hand, a choice is demanded. Making the choice is ipso facto the proof of Dante's poetic vision of the world. Here are some selected passages from the text. Decide where you stand.

\begin{quotex}
Nature was a compound of ideal purposes and inert matter. Life was a conflict between sin and grace. The environment was a battleground between a host of angels and a legion of demons. The better and the worse had actually becomes the sole principles of understanding. 

The highest good became God, the creator of the world. The various stages or elements of perfection became persons in the Godhead, or angelic intelligences, or aerial demons, or lower types of the animal soul. Evil was identified with matter. The various stages of imperfection were ascribed to the grossness of various bodies, which weighted and smothered the spark of divinity that animated them. This spark, however, might be released; then it would fly up again to its parent fire and a soul would be saved. 

Moral values came to be regarded as forces working in nature. But if they worked in nature, which was a compound of evil matter and perfect form, they must exist outside. … The forces that worked in nature were accordingly supernatural virtues, dominations, and powers; each natural thing had its supernatural incubus, a guardian angel, or a devil that possessed it. The supernatural—that is, something moral or ideal regarded as a power and an existence—was all about us. Everything in the world was an effect of something beyond the world; everything in life was a step to something beyond life. 
\end{quotex}

\flrightit{Posted on 2018-03-18 by Cologero }
